\iftestbuild % TEST-ONLY-BEGIN
\setcounter{section}{4}\setcounter{theorem}{70}
\fi % TEST-ONLY-END
% Source: book pp. 175--180 (PDF 189--194); no solutions.
% \axtag and \exnote: see the top of 5C.tex.
\DeclareRobustCommand\axtag[1]{\refstepcounter{theorem}\tag{\thetheorem}\label{#1}}
\section{Commuting Operators}

\begin{definition}[Commute]\label{ax:5.71}\leavevmode
\begin{itemize}
\item Two operators $S$ and $T$ on the same vector space \emph{commute} if
  $ST=TS$.
\item Two square matrices $A$ and $B$ of the same size \emph{commute} if
  $AB=BA$.
\end{itemize}
\end{definition}

For example, if $I$ is the identity operator on $V$ and $\lambda\in\F$, then
$\lambda I$ commutes with every operator on $V$.

As another example, if $T$ is an operator then $T^2$ and $T^3$ commute. More
generally, if $p,q\in\Poly(\F)$, then $p(T)$ and $q(T)$ commute
[see \ref{ax:5.17}(b)].

\begin{example}[Partial differentiation operators commute]\label{ax:5.72}
Suppose $m$ is a nonnegative integer. Let $\Poly_m(\C^2,\C)$ denote the complex
vector space of polynomials (with coefficients in $\C$) in two variables and of
degree at most $m$, with the usual operations of addition and scalar
multiplication of $\C$-valued functions. Thus the elements of $\Poly_m(\C^2,\C)$
are functions $p$ from $\C^2$ to $\C$ of the form
\[ p(w,z) = \sum_{j+k\le m} a_{j,k}w^jz^k, \axtag{ax:5.73} \]
where the indices $j$ and $k$ take on all nonnegative integer values such that
$j+k\le m$, each $a_{j,k}$ is in $\C$, and $w^jz^k$ denotes the function on
$\C^2$ defined by $(w,z)\mapsto w^jz^k$.

Define operators $D_w,D_z\in\Lin\bigl(\Poly_m(\C^2,\C)\bigr)$ by
\[ D_wp = \frac{\partial p}{\partial w} = \sum_{j+k\le m} ja_{j,k}w^{j-1}z^k
   \quad\text{and}\quad
   D_zp = \frac{\partial p}{\partial z} = \sum_{j+k\le m} ka_{j,k}w^jz^{k-1}, \]
where $p$ is as in \ref{ax:5.73}. The operators $D_w$ and $D_z$ are called
partial differentiation operators because each of these operators
differentiates with respect to one of the variables while pretending that the
other variable is a constant.

The operators $D_w$ and $D_z$ commute because if $p$ is as in \ref{ax:5.73},
then
\[ (D_wD_z)p = \sum_{j+k\le m} jka_{j,k}w^{j-1}z^{k-1} = (D_zD_w)p. \]

The equation $D_wD_z=D_zD_w$ on $\Poly_m(\C^2,\C)$ illustrates a more general
result that the order of partial differentiation does not matter for nice
functions.
\end{example}

\marginnote{All 214,358,881 (which equals $11^8$) ordered pairs of the 2-by-2
matrices under consideration were checked by a computer to discover that only
674,609 of these ordered pairs of matrices commute.}%
Commuting matrices are unusual. For example, there are 214,358,881 ordered
pairs of 2-by-2 matrices all of whose entries are integers in the interval
$[-5,5]$. Only about $0.3\%$ of these ordered pairs of matrices commute.

The next result shows that two operators commute if and only if their matrices
(with respect to the same basis) commute.

\begin{result}[Commuting operators correspond to commuting matrices]\label{ax:5.74}
Suppose $S,T\in\Lin(V)$ and $v_1,\dots,v_n$ is a basis of $V$. Then $S$ and $T$
commute if and only if $\Mat\bigl(S,(v_1,\allowbreak\dots,v_n)\bigr)$ and
$\Mat\bigl(T,(v_1,\dots,v_n)\bigr)$ commute.
\end{result}

\begin{proof}
We have
\begin{align*}
S\text{ and }T\text{ commute} &\iff ST=TS\\
  &\iff \Mat(ST)=\Mat(TS)\\
  &\iff \Mat(S)\Mat(T)=\Mat(T)\Mat(S)\\
  &\iff \Mat(S)\text{ and }\Mat(T)\text{ commute},
\end{align*}
as desired.
\end{proof}

The next result shows that if two operators commute, then every eigenspace for
one operator is invariant under the other operator. This result, which we will
use several times, is one of the main reasons why a pair of commuting operators
behaves better than a pair of operators that does not commute.

\begin{result}[Eigenspace is invariant under commuting operator]\label{ax:5.75}
Suppose $S,T\in\Lin(V)$ commute and $\lambda\in\F$. Then $E(\lambda,S)$ is
invariant under $T$.
\end{result}

\begin{proof}
Suppose $v\in E(\lambda,S)$. Then
\[ S(Tv) = (ST)v = (TS)v = T(Sv) = T(\lambda v) = \lambda Tv. \]
The equation above shows that $Tv\in E(\lambda,S)$. Thus $E(\lambda,S)$ is
invariant under $T$.
\end{proof}

Suppose we have two operators, each of which is diagonalizable. If we want to
do computations involving both operators (for example, involving their sum),
then we want the two operators to be diagonalizable by the same basis, which
according to the next result is possible when the two operators commute.

\begin{result}[Simultaneous diagonalizability $\iff$ commutativity]\label{ax:5.76}
Two diagonalizable operators on the same vector space have diagonal matrices
with respect to the same basis if and only if the two operators commute.
\end{result}

\begin{proof}
First suppose $S,T\in\Lin(V)$ have diagonal matrices with respect to the same
basis. The product of two diagonal matrices of the same size is the diagonal
matrix obtained by multiplying the corresponding elements of the two diagonals.
Thus any two diagonal matrices of the same size commute. Thus $S$ and $T$
commute, by \ref{ax:5.74}.

To prove the implication in the other direction, now suppose that
$S,T\in\Lin(V)$ are diagonalizable operators that commute. Let
$\lambda_1,\dots,\lambda_m$ denote the distinct eigenvalues of $S$. Because $S$
is diagonalizable, \ref{ax:5.55}(c) shows that
\[ V = E(\lambda_1,S)\oplus\cdots\oplus E(\lambda_m,S). \axtag{ax:5.77} \]
For each $k=1,\dots,m$, the subspace $E(\lambda_k,S)$ is invariant under $T$
(by \ref{ax:5.75}). Because $T$ is diagonalizable, \ref{ax:5.65} implies that
$T|_{E(\lambda_k,S)}$ is diagonalizable for each $k$. Hence for each
$k=1,\dots,m$, there is a basis of $E(\lambda_k,S)$ consisting of eigenvectors
of $T$. Putting these bases together gives a basis of $V$ (because of
\ref{ax:5.77}), with each vector in this basis being an eigenvector of both $S$
and $T$. Thus $S$ and $T$ both have diagonal matrices with respect to this
basis, as desired.
\end{proof}

See Exercise~2 for an extension of the result above to more than two operators.

Suppose $V$ is a finite-dimensional nonzero complex vector space. Then every
operator on $V$ has an eigenvector (see \ref{ax:5.19}). The next result shows
that if two operators on $V$ commute, then there is a vector in $V$ that is an
eigenvector for both operators (but the two commuting operators might not have
a common eigenvalue). For an extension of the next result to more than two
operators, see Exercise~9(a).

\begin{result}[Common eigenvector for commuting operators]\label{ax:5.78}
Every pair of commuting operators on a finite-dimensional nonzero complex
vector space has a common eigenvector.
\end{result}

\begin{proof}
Suppose $V$ is a finite-dimensional nonzero complex vector space and
$S,T\in\Lin(V)$ commute. Let $\lambda$ be an eigenvalue of $S$ (\ref{ax:5.19}
tells us that $S$ does indeed have an eigenvalue). Thus $E(\lambda,S)\ne\{0\}$.
Also, $E(\lambda,S)$ is invariant under $T$ (by \ref{ax:5.75}).

Thus $T|_{E(\lambda,S)}$ has an eigenvector (again using \ref{ax:5.19}), which
is an eigenvector for both $S$ and $T$, completing the proof.
\end{proof}

\begin{example}[Common eigenvector for partial differentiation operators]\label{ax:5.79}
Let $\Poly_m(\C^2,\C)$ be as in Example~\ref{ax:5.72} and let
$D_w,D_z\in\Lin\bigl(\Poly_m(\C^2,\C)\bigr)$ be the commuting partial
differentiation operators in that example. As you can verify, $0$ is the only
eigenvalue of each of these operators. Also
\begin{align*}
E(0,D_w) &= \biggl\{\sum_{k=0}^m a_kz^k : a_0,\dots,a_m\in\C\biggr\},\\
E(0,D_z) &= \biggl\{\sum_{j=0}^m c_jw^j : c_0,\dots,c_m\in\C\biggr\}.
\end{align*}
The intersection of these two eigenspaces is the set of common eigenvectors of
the two operators. Because $E(0,D_w)\cap E(0,D_z)$ is the set of constant
functions, we see that $D_w$ and $D_z$ indeed have a common eigenvector, as
promised by \ref{ax:5.78}.
\end{example}

The next result extends \ref{ax:5.47} (the existence of a basis that gives an
upper-triangular matrix) to two commuting operators.

\begin{result}[Commuting operators are simultaneously upper triangularizable]\label{ax:5.80}
Suppose $V$ is a finite-dimensional complex vector space and $S,T$ are
commuting operators on $V$. Then there is a basis of $V$ with respect to which
both $S$ and $T$ have upper-triangular matrices.
\end{result}

\begin{proof}
Let $n=\dim V$. We will use induction on $n$. The desired result holds if $n=1$
because all 1-by-1 matrices are upper triangular. Now suppose $n>1$ and the
desired result holds for all complex vector spaces whose dimension is $n-1$.

Let $v_1$ be any common eigenvector of $S$ and $T$ (using \ref{ax:5.78}).
Hence $Sv_1\in\spn(v_1)$ and $Tv_1\in\spn(v_1)$. Let $W$ be a subspace of $V$
such that
\[ V = \spn(v_1)\oplus W; \]
see \ref{ax:2.33} for the existence of $W$. Define a linear map
$P\colon V\to W$ by
\[ P(av_1+w) = w \]
for each $a\in\C$ and each $w\in W$. Define
$\widehat{S},\widehat{T}\in\Lin(W)$ by
\[ \widehat{S}w = P(Sw) \quad\text{and}\quad \widehat{T}w = P(Tw) \]
for each $w\in W$. To apply our induction hypothesis to $\widehat{S}$ and
$\widehat{T}$, we must first show that these two operators on $W$ commute. To
do this, suppose $w\in W$. Then there exists $a\in\C$ such that
\[ (\widehat{S}\widehat{T})w = \widehat{S}\bigl(P(Tw)\bigr)
   = \widehat{S}(Tw-av_1) = P\bigl(S(Tw-av_1)\bigr) = P\bigl((ST)w\bigr), \]
where the last equality holds because $v_1$ is an eigenvector of $S$ and
$Pv_1=0$. Similarly,
\[ (\widehat{T}\widehat{S})w = P\bigl((TS)w\bigr). \]
Because the operators $S$ and $T$ commute, the last two displayed equations
show that $(\widehat{S}\widehat{T})w=(\widehat{T}\widehat{S})w$. Hence
$\widehat{S}$ and $\widehat{T}$ commute.

Thus we can use our induction hypothesis to state that there exists a basis
$v_2,\allowbreak\dots,v_n$ of $W$ such that $\widehat{S}$ and $\widehat{T}$ both have
upper-triangular matrices with respect to this basis. The list
$v_1,\dots,v_n$ is a basis of $V$.

If $k\in\{2,\dots,n\}$, then there exist $a_k,b_k\in\C$ such that
\[ Sv_k = a_kv_1+\widehat{S}v_k \quad\text{and}\quad
   Tv_k = b_kv_1+\widehat{T}v_k. \]
Because $\widehat{S}$ and $\widehat{T}$ have upper-triangular matrices with
respect to $v_2,\dots,v_n$, we know that
$\widehat{S}v_k\in\spn(v_2,\dots,v_k)$ and
$\widehat{T}v_k\in\spn(v_2,\dots,v_k)$. Hence the equations above imply that
\[ Sv_k\in\spn(v_1,\dots,v_k) \quad\text{and}\quad
   Tv_k\in\spn(v_1,\dots,v_k). \]
Thus $S$ and $T$ have upper-triangular matrices with respect to
$v_1,\dots,v_n$, as desired.
\end{proof}

Exercise~9(b) extends the result above to more than two operators.

In general, it is not possible to determine the eigenvalues of the sum or
product of two operators from the eigenvalues of the two operators. However,
the next result shows that something nice happens when the two operators
commute.

\begin{result}[Eigenvalues of sum and product of commuting operators]\label{ax:5.81}
Suppose $V$ is a finite-dimensional complex vector space and $S,T$ are
commuting operators on $V$. Then
\begin{itemize}
\item every eigenvalue of $S+T$ is an eigenvalue of $S$ plus an eigenvalue of
  $T$,
\item every eigenvalue of $ST$ is an eigenvalue of $S$ times an eigenvalue of
  $T$.
\end{itemize}
\end{result}

\begin{proof}
There is a basis of $V$ with respect to which both $S$ and $T$ have
upper-triangular matrices (by \ref{ax:5.80}). With respect to that basis,
\[ \Mat(S+T) = \Mat(S)+\Mat(T) \quad\text{and}\quad \Mat(ST) = \Mat(S)\Mat(T), \]
as stated in \ref{ax:3.35} and \ref{ax:3.43}.

The definition of matrix addition shows that each entry on the diagonal of
$\Mat(S+T)$ equals the sum of the corresponding entries on the diagonals of
$\Mat(S)$ and $\Mat(T)$. Similarly, because $\Mat(S)$ and $\Mat(T)$ are
upper-triangular matrices, the definition of matrix multiplication shows that
each entry on the diagonal of $\Mat(ST)$ equals the product of the
corresponding entries on the diagonals of $\Mat(S)$ and $\Mat(T)$.
Furthermore, $\Mat(S+T)$ and $\Mat(ST)$ are upper-triangular matrices (see
Exercise~2 in Section~5C).

Every entry on the diagonal of $\Mat(S)$ is an eigenvalue of $S$, and every
entry on the diagonal of $\Mat(T)$ is an eigenvalue of $T$ (by \ref{ax:5.41}).
Every eigenvalue of $S+T$ is on the diagonal of $\Mat(S+T)$, and every
eigenvalue of $ST$ is on the diagonal of $\Mat(ST)$ (these assertions follow
from \ref{ax:5.41}). Putting all this together, we conclude that every
eigenvalue of $S+T$ is an eigenvalue of $S$ plus an eigenvalue of $T$, and
every eigenvalue of $ST$ is an eigenvalue of $S$ times an eigenvalue of $T$.
\end{proof}

% ------------------------------------------------------------------ exercises
\begin{exc}

\begin{exercise}
Give an example of two commuting operators $S,T$ on $\F^4$ such that there is a
subspace of $\F^4$ that is invariant under $S$ but not under $T$ and there is a
subspace of $\F^4$ that is invariant under $T$ but not under $S$.
\end{exercise}

\begin{exercise}
Suppose $\mathcal{E}$ is a subset of $\Lin(V)$ and every element of
$\mathcal{E}$ is diagonalizable. Prove that there exists a basis of $V$ with
respect to which every element of $\mathcal{E}$ has a diagonal matrix if and
only if every pair of elements of $\mathcal{E}$ commutes.

\exnote{This exercise extends \ref{ax:5.76}, which considers the case in which
$\mathcal{E}$ contains only two elements. For this exercise, $\mathcal{E}$ may
contain any number of elements, and $\mathcal{E}$ may even be an infinite set.}
\end{exercise}

\begin{exercise}
Suppose $S,T\in\Lin(V)$ are such that $ST=TS$. Suppose $p\in\Poly(\F)$.
\begin{parts}
\item Prove that $\nul p(S)$ is invariant under $T$.
\item Prove that $\range p(S)$ is invariant under $T$.
\end{parts}
\exnote{See \ref{ax:5.18} for the special case $S=T$.}
\end{exercise}

\begin{exercise}
Prove or give a counterexample: If $A$ is a diagonal matrix and $B$ is an
upper-triangular matrix of the same size as $A$, then $A$ and $B$ commute.
\end{exercise}

\begin{exercise}
Prove that a pair of operators on a finite-dimensional vector space commute if
and only if their dual operators commute.

\exnote{See \ref{ax:3.118} for the definition of the dual of an operator.}
\end{exercise}

\begin{exercise}
Suppose that $V$ is a nonzero finite-dimensional complex vector space and
$S,T\in\Lin(V)$ commute. Prove that there exist $\alpha,\lambda\in\C$ such that
\[ \range(S-\alpha I)+\range(T-\lambda I)\ne V. \]
\end{exercise}

\begin{exercise}
Suppose $V$ is a complex vector space, $S\in\Lin(V)$ is diagonalizable, and
$T\in\Lin(V)$ commutes with $S$. Prove that there is a basis of $V$ such that
$S$ has a diagonal matrix with respect to this basis and $T$ has an
upper-triangular matrix with respect to this basis.
\end{exercise}

\begin{exercise}
Suppose $m=3$ in Example~\ref{ax:5.72} and $D_x,D_y$ are the commuting partial
differentiation operators on $\Poly_3(\R^2)$ from that example. Find a basis of
$\Poly_3(\R^2)$ with respect to which $D_x$ and $D_y$ each have an
upper-triangular matrix.
\end{exercise}

\begin{exercise}
Suppose $V$ is a finite-dimensional nonzero complex vector space. Suppose that
$\mathcal{E}\subseteq\Lin(V)$ is such that $S$ and $T$ commute for all
$S,T\in\mathcal{E}$.
\begin{parts}
\item Prove that there is a vector in $V$ that is an eigenvector for every
  element of $\mathcal{E}$.
\item Prove that there is a basis of $V$ with respect to which every element of
  $\mathcal{E}$ has an upper-triangular matrix.
\end{parts}
\exnote{This exercise extends \ref{ax:5.78} and \ref{ax:5.80}, which consider the
case in which $\mathcal{E}$ contains only two elements. For this exercise,
$\mathcal{E}$ may contain any number of elements, and $\mathcal{E}$ may even be
an infinite set.}
\end{exercise}

\begin{exercise}
Give an example of two commuting operators $S,T$ on a finite-dimensional real
vector space such that $S+T$ has an eigenvalue that does not equal an
eigenvalue of $S$ plus an eigenvalue of $T$ and $ST$ has an eigenvalue that
does not equal an eigenvalue of $S$ times an eigenvalue of $T$.

\exnote{This exercise shows that \ref{ax:5.81} does not hold on real vector
spaces.}
\end{exercise}
\end{exc}
